% This standalone fragment does not overwrite the unique-core discussion.
\section{A sharp rank bound for every pure-target optimizer}
\label{sec:all-optimizer-rank}

The existence of a low-rank optimizer does not constrain the rank of other
points in an optimal face. The following theorem supplies a sharp bound for
every optimizer, including singular prescribed marginals.

For positive integers $m,n$ and $1\le k\le\min(m,n)$ define
\[
 N=m+n-k,\qquad
 R_{m,n,k}=\max\{rs:0\le r\le m,\ 0\le s\le n,\ r+s\le N,
                  \ r,s\text{ integers}\}.
\]
The elementary closed form is
\[
 R_{m,n,k}=
 \begin{cases}
 n(m-k),&m-n\ge k,\\
 m(n-k),&n-m\ge k,\\
 \lfloor(m+n-k)^2/4\rfloor,&|m-n|\le k.
 \end{cases}
\]
The formulas coincide at their shared boundaries.

\begin{theorem}[Sharp rank bound for all optimizers]
\label{thm:all-optimizer-rank}
Fix arbitrary density matrices $A,B$ on $\mathbb C^m,\mathbb C^n$ and a
unit pure target $x$ of Schmidt rank $k$. Every coupling maximizing
$\langle x,\rho x\rangle$ with these actual marginals obeys
\[
                  \boxed{\operatorname{rank}\rho\le1+R_{m,n,k}.}
\]
If the optimal target overlap is zero, the stronger bound
$\operatorname{rank}\rho\le R_{m,n,k}$ holds. No positivity assumption
on the marginal eigenvalues is needed.

For every specified unit target of Schmidt rank $k$, except the trivial
$m=n=k=1$ case, there exist positive definite marginals and an optimizer
of rank exactly $1+R_{m,n,k}$. Thus the universal bound is sharp for every
target Schmidt spectrum, not merely for one conveniently chosen target.
In square dimensions it reads
\[
 \boxed{\operatorname{rank}\rho\le
       1+\left\lfloor{(2d-k)^2\over4}\right\rfloor.}
\]
For a full-Schmidt target this is $1+\lfloor d^2/4\rfloor$.
\end{theorem}

\begin{proof}
Let $t=\langle x,\rho x\rangle$ be the optimal value. Suppose first that
$t>0$. The exact extraction
\[
 w={\rho x\over\sqrt t},\qquad R=\rho-ww^\dagger\ge0
\]
satisfies $\langle x,w\rangle=\sqrt t$ and $Rx=0$. Its residual
marginals are
\[
 A_0=A-\operatorname{Tr}_B(ww^\dagger),\qquad
 B_0=B-\operatorname{Tr}_A(ww^\dagger).
\]
They have the same nonnegative trace $\delta$. If $\delta=0$, $R=0$ and
$\rho$ is pure. Otherwise
$ww^\dagger+(A_0\otimes B_0)/\delta$ is another actual coupling.
Its overlap cannot exceed $t$, while its pure term already has overlap
$t$. Consequently
\[
                     \langle x,A_0\otimes B_0\,x\rangle=0.
\]
The product is positive definite on its support, so
$x\perp\operatorname{ran}A_0\otimes\operatorname{ran}B_0$.

Put $r=\operatorname{rank}A_0$ and $s=\operatorname{rank}B_0$.
In local bases adapted to these supports, the coefficient matrix $X$ of
$x$ has a zero $r\times s$ block. Splitting its rows and columns around
that block gives
\[
                k=\operatorname{rank}X\le(m-r)+(n-s),
                \qquad r+s\le N.
\]
Every positive matrix with residual marginals $A_0,B_0$ is supported on
their product support, so $\operatorname{rank}R\le rs$.
Since $Rx=0$ whereas $\langle x,w\rangle\ne0$, the vector $w$ is not in
$\operatorname{ran}R$. For a sum of positive matrices the range is the
sum of their ranges. Thus the ranks add exactly:
\[
        \operatorname{rank}\rho=1+\operatorname{rank}R
                             \le1+rs\le1+R_{m,n,k}.
\]

If $t=0$, the product coupling $A\otimes B$ must also have zero target
overlap. Hence $x\perp\operatorname{ran}A\otimes\operatorname{ran}B$.
Writing $r=\operatorname{rank}A$, $s=\operatorname{rank}B$, the same
zero-block argument gives $r+s\le N$. Every feasible coupling is
supported on that product space, so its rank is at most $rs\le R_{m,n,k}$.

To evaluate $R_{m,n,k}$, increasing either non-saturated variable cannot
decrease $rs$. An optimum may therefore be chosen with $r+s=N$.
Then $m-k\le r\le m$ and one maximizes the concave quadratic
$r(N-r)$ over this integer interval. The closest permitted integer to
$N/2$ gives the displayed piecewise formula.
\end{proof}

\begin{proof}[Sharpness construction]
Choose maximizing integers $r,s$ with $r+s=N$ and put
$a=m-r$, $b=n-s$. Then $a+b=k$. Outside the trivial one-dimensional
case, $r,s>0$. Write a Schmidt decomposition
\[
 x=\sum_{j=1}^{k}\sqrt{p_j}\,e_j\otimes f_j,
 \qquad p_j>0,\quad\sum_jp_j=1.
\]
Let $U$ contain the first $b$ Schmidt vectors on Alice and all of Alice's
unused Schmidt directions. Let $V$ contain the last $a$ Schmidt vectors
on Bob and all of Bob's unused Schmidt directions. Then
\[
 \dim U=b+(m-k)=r,\qquad \dim V=a+(n-k)=s,
\]
and
\[
 x\in(U\otimes V^\perp)\oplus(U^\perp\otimes V),
 \qquad x\perp U\otimes V.
\]
For any $0<q<1$ define
\[
 \rho_q=q\,xx^\dagger+{1-q\over rs}\,P_U\otimes P_V.
\]
It has rank $1+rs$. Both marginals are positive definite: the residual
is strictly positive on $U$ and $V$, while the target's nonzero Schmidt
coefficients cover their respective complements.

The local-sum operator
\[
          D=P_{U^\perp}\otimes I+I\otimes P_{V^\perp}
\]
is zero on $U\otimes V$, the identity on the two off-diagonal blocks
containing $x$, and twice the identity on
$U^\perp\otimes V^\perp$. Therefore $D\ge xx^\dagger$ and
$\operatorname{Tr}(D\rho_q)=q$. Every state with the same two marginals
has this same local-sum expectation and hence target overlap at most
$q$. The displayed state has target overlap exactly $q$, so it is an
optimizer attaining the bound. When $m=n=k=1$, the only coupling has
rank one, as claimed.
\end{proof}

\paragraph{Relation to prior rank bounds.}
Parthasarathy's general extremal-coupling bound is
$\operatorname{rank}\rho\le\sqrt{m^2+n^2-1}$
(arXiv:quant-ph/0307182). That is a statement about extreme points of the
whole coupling set. The theorem above instead concerns every point of
the exposed face selected by a pure target and tracks the target's Schmidt
rank. It complements the separate existence theorem giving an optimizer
of rank at most $\max\{\operatorname{rank}A,\operatorname{rank}B\}$. A high-rank member of the optimal face need not
be extreme. These distinctions are essential when comparing the bounds.
The precise target-dependent sharp theorem was not located in the bounded
primary-source search; historical priority has not been certified.

\paragraph{Verification scope.}
The proof is exact. The supporting rank-one extraction, positive completion,
and zero-residual-expectation steps are supplied in the finite-matrix Lean
modules. The full support/rank theorem is not claimed to be Lean formalized.
The accompanying exact-arithmetic script verifies the maximizing integer
formula, the sharp states and their local-sum dual certificates for small
rectangular dimensions and all allowed Schmidt ranks.

\subsection{What the two rank statements say in small dimensions}
For full-Schmidt targets on square systems, the existence bound and the
all-optimizer bound are both sharp, but answer different questions.
\begin{center}
\begin{tabular}{ccl}\toprule
Local dimension $d$ & Some optimizer & Every optimizer\\\midrule
2 & rank at most 2 & rank at most 2\\
3 & rank at most 3 & rank at most 3\\
4 & rank at most 4 & rank at most 5\\
5 & rank at most 5 & rank at most 7\\
6 & rank at most 6 & rank at most 10\\
8 & rank at most 8 & rank at most 17\\\bottomrule
\end{tabular}
\end{center}
For positive definite marginals, fixed full-Schmidt targets in dimensions
two and three have a unique optimizer, by the unique-core theorem below.
Dimension four permits both a nine-dimensional optimal face and rank-five
optimizers. Thus a low-rank attaining representative must not be mistaken
for a classification of all attaining states.
